\documentclass[12pt]{article}
\usepackage[utf8]{inputenc}
\usepackage[T1]{fontenc}
\usepackage{hyperref}
\title{Vidunda Language Resources for K-12 Curriculum Development}
\author{Compiled by Claude Sonnet 5 (claude-sonnet-5) for Andr\'as Kornai}
\date{}
\begin{document}
\maketitle

\section{Introduction}
Vidunda (Kividunda, Chividunda; ISO 639-3 vid) is a Bantu language spoken along the north bank of the Ruaha River in Kilosa District, Morogoro Region, Tanzania. Speaker estimates range from about 12,000 to 32,000. Endangered Languages lists it as threatened: children learn it at home, but by Standard 7 most use Swahili with peers. The searches found a small body of documentation and no school materials.

\section{Primary Educational Resources}
None found.

\section{Linguistic and Reference Materials}
Karsten Legère and Peter Mkwan'hembo have documented the language, including a 2020 proverb collection with Swahili translations.
\begin{itemize}
\item \emph{Luhala lwa Wavidunda -- Vidunda Wisdom} (G.38), proverbs, R\"udiger K\"oppe Verlag: \url{https://www.koeppe.de/titel_luhala-lwa-wavidunda-vidunda-wisdom-g-38}
\item Publisher shop page for the same title: \url{https://shop.koeppe.de/en/?p=4089}
\item Article on Vidunda (journal.fi): \url{https://journal.fi/store/article/view/52597/16362}
\item Vidunda language (Wikipedia): \url{https://en.wikipedia.org/wiki/Vidunda_language}
\item Vidunda people (Wikipedia): \url{https://en.wikipedia.org/wiki/Vidunda_people}
\item Endangered Languages entry: \url{https://endangeredlanguages.com/lang/1101?hl=en}
\end{itemize}

\section{Religious and Community Resources}
A Global Recordings Network page for vid exists, which suggests gospel recordings: \url{https://glorecnet-website.k.globalrecordings.net/en/language/vid}. I did not inspect its content.

\section{Digital and Audio Resources}
See the Global Recordings Network page above.

\section{Curriculum Development Context}
Vidunda is not known to be taught in schools, and Swahili dominates by upper primary. The proverb collection and Legère's documentation work are the best starting points for local-language content.
\end{document}